%
% Chapter 3: Data Collection
%
\chapter{Data Collection}
\label{chap:data-collection}
\setlength{\parskip}{1em}

This chapter describes the data sources, benchmark setup, and evaluation metrics used to assess Snowdrop’s performance. The main validation framework is the TPC-H benchmark, which is widely used as an industry standard for measuring analytical query performance in decision-support systems.

\section{TPC-H Benchmark Overview}

The Transaction Processing Performance Council Benchmark H (TPC-H) is a standardized benchmark designed to evaluate database systems on analytical workloads~\cite{tpch2024}. It defines a schema consisting of eight related tables that model a wholesale supplier scenario, along with 22 SQL queries covering a broad range of operations, including scans, aggregations, joins, subqueries, and sorting.

In this study, Snowdrop is evaluated using TPC-H at scale factor~1, which corresponds to approximately 1\,GB of raw data. The dataset is generated with the \texttt{tpch-dbgen} tool, which produces pipe-delimited flat files for each table. These files are then loaded into both Snowdrop and PostgreSQL~16 on the same hardware in order to keep the comparison as fair as possible.

Scale factor~1 was selected for two main reasons. First, it is large enough to reveal meaningful differences between storage and execution architectures. For example, the \texttt{lineitem} table alone contains around 6 million rows across 16 columns. Second, the dataset is still small enough to fit entirely in memory on a single development machine, which helps focus the evaluation on query execution efficiency rather than on raw disk throughput.

\section{Dataset Structure}

Table~\ref{tab:tpch-schema} presents the structure of the TPC-H dataset at scale factor~1.

\begin{table}[H]
\centering
\caption{TPC-H dataset structure at scale factor 1}
\label{tab:tpch-schema}
\renewcommand{\arraystretch}{1.4}
\begin{tabular}{|l|r|r|l|}
\hline
\textbf{Table} & \textbf{Rows} & \textbf{Columns} & \textbf{Indexed columns (Snowdrop)} \\
\hline
lineitem  & 6,000,000 & 16 & l\_shipdate, l\_orderkey, l\_partkey \\
\hline
orders    & 1,500,000 & 9  & o\_orderdate, o\_custkey \\
\hline
customer  & 150,000   & 8  & c\_custkey \\
\hline
partsupp  & 800,000   & 5  & --- \\
\hline
part      & 200,000   & 9  & --- \\
\hline
supplier  & 10,000    & 7  & --- \\
\hline
nation    & 25        & 4  & --- \\
\hline
region    & 5         & 3  & --- \\
\hline
\end{tabular}
\end{table}

The three largest tables --- \texttt{lineitem}, \texttt{orders}, and \texttt{customer} --- make up more than 99\% of the total number of rows and therefore become the main focus of the performance evaluation. In Snowdrop, secondary indexes are created on columns that appear in the WHERE clauses of the selected TPC-H queries. In the current architecture these indexed values are stored in RocksDB alongside the durable row data, while Neo4j retains only schema metadata describing which columns are indexed.

\section{Target Queries}

Four TPC-H queries were chosen as the main benchmark targets, with each query intended to test a different aspect of Snowdrop’s architecture. Table~\ref{tab:target-queries} summarizes these queries and the specific capabilities they help validate.

\begin{table}[H]
\centering
\caption{Target TPC-H queries and validated capabilities}
\label{tab:target-queries}
\renewcommand{\arraystretch}{1.4}
\small
\begin{tabular}{|l|p{4cm}|p{5.5cm}|}
\hline
\textbf{Query} & \textbf{What it tests} & \textbf{Snowdrop capability validated} \\
\hline
Q6  & Date range scan, SUM over 6M rows & Columnar cache, SIMD filtering, aggregation \\
\hline
Q1  & GROUP BY with 7 aggregates, date range & Full aggregate pipeline, ORDER BY \\
\hline
Q14 & 2-table JOIN, date range, string match & JOIN support, filter buffer predicates \\
\hline
Q3  & 3-table JOIN with GROUP BY, ORDER BY & Multi-table JOIN, aggregates, sorting \\
\hline
\end{tabular}
\end{table}

Q6 is treated as the primary validation query because it exercises Snowdrop’s main distinguishing feature: the columnar cache combined with SIMD-accelerated predicate evaluation. The query applies five predicates to the \texttt{lineitem} table and computes a SUM aggregate over the matching rows. Q1 extends this scenario to a multi-column GROUP BY with seven aggregate functions. Q14 and Q3 introduce multi-table JOIN operations, which stress a different part of the system and make Snowdrop’s current limitations in optimizer quality and multi-table execution more visible.

\section{Performance Metrics}

Each query is executed under two different conditions in order to separate initialization costs from steady-state execution behavior.

\textbf{Cold cache:} The orchestrator is restarted before each run, which forces Snowdrop to rebuild the columnar cache from RocksDB during the first scan. This captures the combined cost of cache construction and query execution.

\textbf{Warm cache:} The query is executed after a prior run has already populated the columnar cache. This allows the evaluation to focus on steady-state execution performance and better reflects the latency that repeated workloads would observe in practice.

Under each condition, every query is executed 10 times. The following metrics are collected:

\begin{itemize}
  \item \textbf{p50 latency} --- the median response time across 10 runs
  \item \textbf{p95 latency} --- the 95th percentile response time
  \item \textbf{p99 latency} --- the 99th percentile response time, which serves as Snowdrop’s primary optimization target
  \item \textbf{Mean latency} --- the arithmetic mean across all runs
\end{itemize}

\section{Hardware and Environment}

All benchmarks are run on the same hardware for both Snowdrop and PostgreSQL~16 in order to preserve comparability. The test environment includes the following components:

\begin{itemize}
  \item Docker Compose orchestration with dedicated containers for Neo4j~5, PostgreSQL~16, and the Rust orchestrator, with optional Kvrocks preserved only for architectural comparison
  \item Neo4j configured with a 2\,GB initial heap and a 4\,GB maximum heap, using the ZGC garbage collector (\texttt{-XX:+UseZGC})
  \item RocksDB with default level compaction and Bloom filters enabled
  \item PostgreSQL~16 running with its default configuration and without additional tuning, representing an out-of-the-box baseline
\end{itemize}

The choice to compare against an untuned PostgreSQL~16 instance is intentional. The purpose of the evaluation is not to claim that Snowdrop outperforms a production-optimized PostgreSQL deployment in all cases, but rather to show that the proposed hybrid architecture can reach competitive performance on selected analytical queries while preserving SQL compatibility. A more heavily tuned PostgreSQL instance, with workload-specific indexes, adjusted shared buffers, and parallel query settings, would likely reduce the performance gap on several benchmarks.

\section{Profiling Tools}

In addition to raw latency measurements, deeper performance analysis is carried out with two profiling tools.

\textbf{tokio-console} is an async runtime profiler for Tokio-based Rust applications. It provides real-time visibility into task scheduling, polling duration, and resource usage inside the orchestrator. This makes it possible to determine whether query latency is primarily caused by computation, I/O waiting, or scheduling overhead.

\textbf{cargo flamegraph} generates CPU flamegraphs from the orchestrator binary and shows how execution time is distributed across function calls. This is especially useful for detecting unexpected bottlenecks. For example, the earlier finding that Kvrocks TCP round-trips accounted for about 70\% of total query time was identified through flamegraph analysis~\cite{luo2019}.

Both tools are used starting from Milestone~4, once the query execution pipeline is stable enough to produce representative performance profiles.
